\textbf{Phần 1. Trắc nghiệm lựa chọn}
\setcounterpageref{ex}{0}

\begin{ex} %[3D2V11]
Tính giá trị của biểu thức: $C = 8 \times 9 - 56 : 8 + 24$.
\choice
{\True $89$}
{$85$}
{$91$}
{$87$}
\loigiai{
Thực hiện nhân chia trước: $C = 72 - 7 + 24 = 65 + 24 = 89$.
\textbf{Đáp án của câu hỏi là A}
}
\end{ex}

\begin{ex} %[3D2V11]
Một con bạch tuộc có $8$ cái tua. Hỏi một đàn gồm $9$ con bạch tuộc có tất cả bao nhiêu cái tua?
\choice
{$64$ cái tua}
{$80$ cái tua}
{\True $72$ cái tua}
{$68$ cái tua}
\loigiai{
Số cái tua là: $9 \times 8 = 72$ (cái tua).
\textbf{Đáp án của câu hỏi là C}
}
\end{ex}

\begin{ex} %[3D2V11]
Tìm $x$ biết: $x : 8 + 35 = 43$.
\choice
{$56$}
{\True $64$}
{$72$}
{$48$}
\loigiai{
Ta có: $x : 8 = 43 - 35 = 8 \Rightarrow x = 8 \times 8 = 64$.
\textbf{Đáp án của câu hỏi là B}
}
\end{ex}

\begin{ex} %[3D2V11]
Trong một phép chia cho $8$, thương bằng $7$ và số dư là số dư lớn nhất có thể có. Số bị chia của phép chia đó là
\choice
{$56$}
{$61$}
{$62$}
{\True $63$}
\loigiai{
Vì số chia là 8 nên số dư lớn nhất là 7.\
Số bị chia là: $8 \times 7 + 7 = 56 + 7 = 63$.
\textbf{Đáp án của câu hỏi là D}
}
\end{ex}

\begin{ex} %[3D2V11]
Tính nhanh: $8 \times 2 + 8 \times 3 + 8 \times 5$.
\choice
{\True $80$}
{$72$}
{$88$}
{$96$}
\loigiai{
Đặt thừa số chung: $8 \times (2 + 3 + 5) = 8 \times 10 = 80$.
\textbf{Đáp án của câu hỏi là A}
}
\end{ex}

\begin{ex} %[3D2V11]
Một cuộn dây cáp dài $64\text{ m}$. Người ta cắt thành $8$ đoạn dây dài bằng nhau. Hỏi $5$ đoạn dây như thế dài bao nhiêu mét?
\choice
{$35\text{ m}$}
{$45\text{ m}$}
{\True $40\text{ m}$}
{$48\text{ m}$}
\loigiai{
Mỗi đoạn dây dài: $64 : 8 = 8\text{ (m)}$.\
5 đoạn dây dài: $8 \times 5 = 40\text{ (m)}$.
\textbf{Đáp án của câu hỏi là C}
}
\end{ex}

\begin{ex} %[3D2V11]
Khối lập phương có đúng 8 đỉnh. Hỏi 6 khối lập phương như thế có tất cả bao nhiêu đỉnh?
\choice
{$40$ đỉnh}
{\True $48$ đỉnh}
{$54$ đỉnh}
{$56$ đỉnh}
\loigiai{
Số đỉnh của 6 khối lập phương là: $8 \times 6 = 48$ (đỉnh).
\textbf{Đáp án của câu hỏi là B}
}
\end{ex}

\begin{ex} %[3D2V11]
Tìm $x$ biết: $(x - 24) \times 8 = 48$.
\choice
{$28$}
{$32$}
{$35$}
{\True $30$}
\loigiai{
Ta có: $x - 24 = 48 : 8 = 6 \Rightarrow x = 6 + 24 = 30$.
\textbf{Đáp án của câu hỏi là D}
}
\end{ex}

\begin{ex} %[3D2V11]
Một cửa hàng có $72\text{ lít}$ mật ong chia đều vào 8 can. Cửa hàng đã bán được 3 can mật ong. Hỏi cửa hàng còn lại bao nhiêu lít mật ong?
\choice
{\True $45\text{ lít}$}
{$40\text{ lít}$}
{$50\text{ lít}$}
{$48\text{ lít}$}
\loigiai{
Mỗi can chứa: $72 : 8 = 9\text{ lít}$.\
Số can còn lại: $8 - 3 = 5$ can.\
Số lít còn lại: $5 \times 9 = 45\text{ lít}$.
\textbf{Đáp án của câu hỏi là A}
}
\end{ex}

\begin{ex} %[3D2V11]
Cho dãy số: $8; 16; 24; 32; \dots$ Số hạng thứ 9 của dãy số trên là
\choice
{$64$}
{$80$}
{\True $72$}
{$68$}
\loigiai{
Số hạng thứ 9 là: $8 \times 9 = 72$.
\textbf{Đáp án của câu hỏi là C}
}
\end{ex}

\begin{ex} %[3D2V11]
Một người thợ làm việc $8$ giờ mỗi ngày và làm việc trong 6 ngày một tuần. Hỏi trong 2 tuần người đó làm việc tất cả bao nhiêu giờ?
\choice
{$88$ giờ}
{\True $96$ giờ}
{$92$ giờ}
{$100$ giờ}
\loigiai{
Mỗi tuần người đó làm: $8 \times 6 = 48$ giờ.\
Trong 2 tuần làm: $48 \times 2 = 96$ giờ.
\textbf{Đáp án của câu hỏi là B}
}
\end{ex}

\begin{ex} %[3D2V11]
Điền dấu thích hợp: $8 \times 7 - 16 \dots 40 : 8 \times 8$.
\choice
{$>$}
{$<$}
{Không so sánh được}
{\True $=$}
\loigiai{
Vế trái: $8 \times 7 - 16 = 56 - 16 = 40$.\
Vế phải: $40 : 8 \times 8 = 5 \times 8 = 40$.\
Do đó điền dấu $=$.
\textbf{Đáp án của câu hỏi là D}
}
\end{ex}

\textbf{Phần 2. Trắc nghiệm đúng sai}
\setcounterpageref{ex}{0}

\begin{ex} %[3D2V11]
Xét tính đúng sai của các nhận định sau về phép nhân và chia với 8:
\choiceTF
{\True Tích của số 8 với một số tự nhiên bất kì luôn là số chẵn.}
{Số chia hết cho 8 thì không chia hết cho 4.}
{\True Trong phép chia cho 8, số dư lớn nhất có thể có là 7.}
{\True Giá trị của $8 \times 5$ bằng một nửa của $8 \times 10$.}
\loigiai{
\begin{itemize}
    \item Ý a) Đúng: Vì 8 là số chẵn.
    \item Ý b) Sai: Mọi số chia hết cho 8 đều chia hết cho 4 vì $8 = 4 \times 2$.
    \item Ý c) Đúng: Số dư luôn bé hơn số chia.
    \item Ý d) Đúng: Vì $40$ bằng một nửa của $80$.
\end{itemize}
\begin{itemize}
    \item a) \textbf{Đúng}.
    \item b) \textbf{Sai}.
    \item c) \textbf{Đúng}.
    \item d) \textbf{Đúng}.
\end{itemize}
}
\end{ex}

\begin{ex} %[3D2V11]
Một thùng kẹo có 48 cái kẹo chia đều vào 8 hộp nhỏ:
\choiceTF
{\True Mỗi hộp nhỏ có 6 cái kẹo.}
{\True Lấy ra 3 hộp kẹo thì được tất cả 18 cái kẹo.}
{Số kẹo trong 5 hộp nhỏ là 32 cái kẹo.}
{\True Tổng số kẹo gấp 8 lần số kẹo của một hộp nhỏ.}
\loigiai{
\begin{itemize}
    \item Ý a) Đúng: $48 : 8 = 6$ cái kẹo.
    \item Ý b) Đúng: $3 \times 6 = 18$ cái kẹo.
    \item Ý c) Sai: $5 \times 6 = 30$ cái kẹo.
    \item Ý d) Đúng: $48 : 6 = 8$ lần.
\end{itemize}
\begin{itemize}
    \item a) \textbf{Đúng}.
    \item b) \textbf{Đúng}.
    \item c) \textbf{Sai}.
    \item d) \textbf{Đúng}.
\end{itemize}
}
\end{ex}

\textbf{Phần 3. Trả lời ngắn}
\setcounterpageref{ex}{0}

\begin{ex} %[3D2V11]
Tính giá trị biểu thức: $8 \times 6 + 32 : 8$.
\shortans{52}
\loigiai{
\textbf{Đáp số:} $52$\\
$8 \times 6 + 32 : 8 = 48 + 4 = 52$.
}
\end{ex}

\begin{ex} %[3D2V11]
Một con nhện có 8 cái chân. Hỏi 7 con nhện có tất cả bao nhiêu cái chân?
\shortans{56}
\loigiai{
\textbf{Đáp số:} $56$\\
$7 \times 8 = 56$ cái chân.
}
\end{ex}

\begin{ex} %[3D2V11]
Tìm số bị chia trong phép chia cho 8, biết thương là 6 và số dư là 3.
\shortans{51}
\loigiai{
\textbf{Đáp số:} $51$\\
$8 \times 6 + 3 = 48 + 3 = 51$.
}
\end{ex}

\begin{ex} %[3D2V11]
Tìm $x$ biết: $x : 8 = 15 - 9$.
\shortans{48}
\loigiai{
\textbf{Đáp số:} $48$\\
$x : 8 = 6 \Rightarrow x = 6 \times 8 = 48$.
}
\end{ex}

\begin{ex} %[3D2V11]
Một lớp học có 32 học sinh chia đều thành 8 tổ. Hỏi 5 tổ có bao nhiêu học sinh?
\shortans{20}
\loigiai{
\textbf{Đáp số:} $20$\\
Mỗi tổ có: $32 : 8 = 4$ bạn. 5 tổ có: $5 \times 4 = 20$ bạn.
}
\end{ex}

\begin{ex} %[3D2V11]
Tính nhanh: $8 \times 1 + 8 \times 2 + 8 \times 3 + 8 \times 4$.
\shortans{80}
\loigiai{
\textbf{Đáp số:} $80$\\
$8 \times (1 + 2 + 3 + 4) = 8 \times 10 = 80$.
}
\end{ex}
